\chapter{Eficiência energética e sustentabilidade}
\label{cap:eficiencia}

\begin{objetivos}
Ao final deste capítulo você deve ser capaz de:
\begin{itemize}
  \item identificar, na lógica de controle, as decisões que alteram o consumo de energia;
  \item definir indicadores de consumo comparáveis entre períodos, linhas e produtos;
  \item escolher entre partida direta, partida suave e acionamento por variador de frequência;
  \item aplicar controle de demanda sem comprometer a segurança do processo;
  \item considerar o ciclo de vida do equipamento, da escolha ao descarte.
\end{itemize}
\end{objetivos}

\section{A lógica também consome energia}

Eficiência energética em automação costuma ser tratada como assunto de equipamento --- trocar motor,
instalar variador, colocar painel solar. Boa parte do resultado, porém, está em \emph{quando} e
\emph{como} a máquina opera, e isso é decidido na lógica.

Três exemplos do dia a dia em que o programa decide o consumo:

\begin{itemize}
  \item \textbf{operação a vazio}: a esteira continua ligada depois de a última peça passar. É o
        consumo mais fácil de eliminar e o mais frequentemente esquecido;
  \item \textbf{excesso de tempo de espera}: um temporizador longo demais mantém aquecedor, bomba ou
        exaustor ligados sem necessidade;
  \item \textbf{pico de partida desnecessário}: todos os motores partem juntos, elevando a demanda
        instantânea --- e a fatura --- sem ganho de produção.
\end{itemize}

\begin{atencao}
Eficiência nunca compete com segurança. Desligar automaticamente um sistema de exaustão, uma
pressurização, uma refrigeração de painel ou qualquer função de proteção para economizar energia é
falha grave de projeto. Antes de automatizar qualquer desligamento, confirme por escrito quais
funções \emph{não} podem ser interrompidas, ouvindo quem opera e quem mantém --- e registre essa lista
na documentação (Capítulo~\ref{cap:documentacao}).
\end{atencao}

\section{Medir antes de agir}

Sem medição, toda ação de eficiência é opinião. O mínimo necessário:

\begin{enumerate}
  \item \textbf{medidor de energia} no ponto que se quer avaliar: quadro de uma linha, de um grupo de
        máquinas ou de um equipamento que consome muito;
  \item \textbf{leitura por rede} quando o medidor oferece comunicação, ou por \textbf{pulso} em
        entrada digital, com o cuidado de escolher a resolução adequada à frequência do pulso
        (Capítulo~\ref{cap:es});
  \item \textbf{energia por unidade produzida}, que é o indicador comparável. Quilowatt-hora por
        tonelada, por peça ou por metro produzido permite comparar turnos e produtos diferentes, o
        que o consumo absoluto não permite;
  \item \textbf{contexto}: produto, velocidade, turno e temperatura ambiente. Comparar consumo de
        dois períodos com produtos diferentes e concluir que a eficiência mudou é erro clássico
        (Capítulo~\ref{cap:dados}).
\end{enumerate}

\begin{figure}[H]
  \centering
  \includegraphics[width=0.42\linewidth]{23-alavancas-consumo.pdf}
  \caption{O ciclo da eficiência: medir, agir na lógica e verificar pelo indicador.}
  \label{fig:alavancas-consumo}
\end{figure}

A Figura~\ref{fig:alavancas-consumo} tem um retorno explícito: a seta pontilhada significa que a
ação precisa ser reavaliada pelo indicador. Alavanca instalada sem verificação é despesa, não
resultado.

Com notação de grandezas: a energia $E$ em quilowatt-hora é o produto da potência média $\bar{P}$ em
quilowatt pelo intervalo de tempo $\Delta t$ em horas,

\[
E = \bar{P} \cdot \Delta t ,
\]

e o indicador de eficiência é a razão entre a energia e a produção $Q$ no mesmo período:

\[
e = \frac{E}{Q} \quad [\text{kWh por unidade produzida}] .
\]

É essa razão --- e não $E$ sozinho --- que deve entrar em meta. Produzir mais consumindo
proporcionalmente mais não é eficiência.

\index{eficiência energética}\index{kWh por unidade produzida}

\section{Alavancas na lógica}

\begin{table}[H]
  \centering
  \caption{Alavancas de eficiência que se implementam na lógica de controle.}
  \label{tab:alavancas-eficiencia}
  \begin{tabular}{p{3.7cm} p{5.6cm} p{5.0cm}}
    \toprule
    Alavanca & O que fazer no programa & Indicador afetado \\
    \midrule
    Desligamento automático por inatividade
      & Temporizar a ausência de peça ou de demanda e desligar o que pode ser desligado: esteira,
        exaustão, iluminação de área, aquecimento de espera
      & Consumo da linha no turno \\
    \addlinespace
    Fim da operação a vazio
      & Detectar ``sem carga'' e reduzir rotação ou parar, em vez de manter o equipamento girando
      & Energia por peça produzida \\
    \addlinespace
    \emph{Setpoint} ótimo
      & Revisar referências que ficaram com folga grande (temperatura, pressão, vazão) e ajustar ao
        necessário, com confirmação de que a qualidade se mantém
      & Consumo específico do processo \\
    \addlinespace
    Ajuste por demanda
      & Reduzir rotação de bomba ou ventilador quando a demanda cai, em vez de manter rotação
        constante com estrangulamento
      & Energia do sistema de bombeamento ou ventilação \\
    \addlinespace
    Redução de tempo de espera
      & Encurtar temporizações de aquecimento, estabilização e transição, dentro do que o processo
        exige
      & Energia por lote \\
    \addlinespace
    Sequenciamento de partida
      & Escalonar a partida de motores para reduzir a demanda instantânea
      & Demanda contratada e multa por ultrapassagem \\
    \addlinespace
    Recuperação de calor e reúso
      & Usar energia que já está no processo: reaproveitar calor de compressor, de exaustão ou de
        resfriamento
      & Consumo total da planta \\
    \bottomrule
  \end{tabular}
\end{table}

\begin{atencao}
Ajuste de \emph{setpoint} é a alavanca com melhor retorno e maior risco. Reduzir a temperatura de um
banho, a pressão de uma linha de ar ou a vazão de resfriamento pode afetar qualidade, tempo de ciclo
ou segurança de processo. Toda alteração de referência precisa de teste documentado e aprovação de
quem responde pelo produto e pela instalação --- nunca de ajuste anônimo em tempo de execução.
\end{atencao}

\section{Acionamento de motores}

A escolha do modo de partida tem efeito sobre o consumo, sobre a rede elétrica e sobre a vida útil do
equipamento:

\begin{itemize}
  \item \textbf{partida direta}: simples e barata. A corrente de partida é várias vezes a nominal, o
        que produz queda de tensão, esforço mecânico e pico de demanda. Adequada a motores pequenos e
        a partidas pouco frequentes;
  \item \textbf{partida estrela-triângulo}: reduz a corrente de partida, com degrau de conjugado e
        limitação de carga no momento da transição. Serve a cargas que partem descarregadas;
  \item \textbf{partida suave} (\emph{soft-starter}): controla a tensão aplicada na partida, sem
        controlar a rotação em regime. É solução de compromisso para reduzir esforço mecânico e pico
        de corrente;
  \item \textbf{variador de frequência}: controla a rotação em regime. É o único dos quatro que
        permite eficiência energética em carga variável, porque atua no ponto de operação, e não
        apenas na partida.
\end{itemize}

Em bombas centrífugas e ventiladores, o ponto de operação define o consumo. Reduzir a vazão fechando
uma válvula ou um \emph{damper} aumenta a perda de carga e mantém o motor consumindo quase o mesmo;
reduzir a rotação, por outro lado, reduz o consumo de forma muito mais acentuada --- em sistemas
centrífugos a potência varia de forma aproximadamente cúbica com a rotação, o que faz da
regulagem por rotação a alavanca mais forte disponível em bombeamento e ventilação.

Isso não significa que variador é sempre a resposta: em carga constante, o inversor acrescenta perdas
próprias e complexidade sem ganho. A pergunta correta é ``a carga varia?''.

\index{variador de frequência}\index{partida de motor}

\section{Controle de demanda}

A fatura de energia industrial depende tanto da \emph{energia} consumida quanto da \emph{demanda}
--- a potência média no intervalo de medição --- e multa a ultrapassagem do valor contratado. Dois
mecanismos reagem a isso:

\begin{enumerate}
  \item \textbf{limitação por prioridade}: quando a demanda medida se aproxima do limite, o sistema
        desliga cargas de menor prioridade por tempo determinado. É preciso que a lista de prioridades
        seja explícita e que as funções críticas nunca estejam nela;
  \item \textbf{gerenciamento por horário}: deslocar operações com folga --- carregamento de baterias,
        limpeza, teste de equipamento --- para fora do horário de ponta.
\end{enumerate}

A limitação de demanda é uma intervenção \emph{na operação}, e por isso precisa de três cuidados:
lista de cargas bloqueadas revisada e aprovada, alarme claro para o operador quando a limitação
atuar, e registro de cada atuação --- sem isso, a produção perde uma máquina e a causa fica
desconhecida.

\section{Ciclo de vida e descarte}

Eficiência e sustentabilidade também se decidem antes da instalação:

\begin{itemize}
  \item \textbf{vida útil e disponibilidade de reposição}: equipamento com fim de vida anunciado
        força substituição no pior momento, e a substituição emergencial consome mais recurso que a
        planejada (Capítulo~\ref{cap:manutencao});
  \item \textbf{atualização sem troca total}: parte dos casos de obsolescência se resolve com
        atualização de \emph{firmware}, troca de módulo ou de HMI, preservando o que ainda funciona;
  \item \textbf{sobressalentes críticos} em estoque, com regra de descarte por validade ---
        especialmente baterias de retenção e filtros;
  \item \textbf{descarte correto}: eletroeletrônicos, baterias e componentes não seguem o lixo comum;
        seguem logística reversa, conforme a política de resíduos adotada pela instituição ou pela
        empresa;
  \item \textbf{eficiência do próprio sistema de automação}: fontes sobredimensionadas, painéis mal
        ventilados e iluminação de painel permanente são consumos pequenos que se somam em uma planta
        inteira.
\end{itemize}

\begin{pratica}
Escolha um equipamento de consumo significativo na sua planta ou bancada e aplique o ciclo da
Figura~\ref{fig:alavancas-consumo}:
(a) instale ou identifique a medição e calcule a energia por unidade produzida no período;
(b) escolha duas alavancas da Tabela~\ref{tab:alavancas-eficiencia} que se apliquem ao equipamento;
(c) escreva a alteração de lógica necessária para cada uma, em texto estruturado;
(d) defina o que \emph{não} pode ser desligado e registre a justificativa;
(e) proponha o prazo de reavaliação do indicador e o critério de sucesso.
\end{pratica}

\begin{exercicios}
\begin{enumerate}[label=\alph*)]
  \item Dê três exemplos de consumo causado por decisão de programação, e não por equipamento.
  \item Por que o indicador ``kWh por tonelada'' é preferível ao consumo absoluto em kWh?
  \item Escreva a expressão da energia a partir de potência média e intervalo de tempo, e explique
        cada grandeza com a unidade.
  \item Uma linha consumiu 4.800\,kWh para produzir 60\,t em um turno e 5.200\,kWh para produzir 65\,t
        no turno seguinte. Calcule o indicador nos dois turnos e diga o que melhorou.
  \item Explique por que a limitação de demanda exige lista de prioridades aprovada e registro de
        atuação.
  \item Compare partida direta, estrela-triângulo, partida suave e acionamento por variador quanto a
        corrente de partida e controle de regime.
  \item Por que, em bombeamento, reduzir a rotação economiza mais que fechar uma válvula? Relacione
        com o ponto de operação.
  \item Um colega propõe instalar variador em um motor que trabalha sempre em carga nominal. Faça a
        análise técnica da proposta.
  \item Cite dois perigos de reduzir \emph{setpoint} de temperatura em um processo, e diga quem deve
        aprovar a alteração.
  \item Selecione duas alavancas da Tabela~\ref{tab:alavancas-eficiencia} e escreva o trecho de lógica
        necessário para cada uma.
  \item Liste três funções que \emph{nunca} podem entrar em um esquema de desligamento automático, e
        justifique cada uma.
  \item Descreva duas medidas de sustentabilidade ligadas ao ciclo de vida do controlador, e diga em
        que momento do projeto cada uma deve ser decidida.
  \item Um painel elétrico tem as fontes sobredimensionadas e a ventilação desligada nos fins de
        semana. Explique, com base neste capítulo, se a decisão é correta.
\end{enumerate}
\end{exercicios}
